% =====================================================================
% INSERT FOR PAPER IV  (Arithmetic Expression Geometry IV)
%
% Source file:  paper-4/sections/04b-transport-density.tex
% Integrated 2026-10-20 as the last subsection of Section 4, immediately after
%   Subsection 4.7 ("Non-exact projective transport and cross-ratio holonomy")
%   and its Example 4.13, via
%   % =====================================================================
% INSERT FOR PAPER IV  (Arithmetic Expression Geometry IV)
%
% Source file:  paper-4/sections/04b-transport-density.tex
% Integrated 2026-10-20 as the last subsection of Section 4, immediately after
%   Subsection 4.7 ("Non-exact projective transport and cross-ratio holonomy")
%   and its Example 4.13, via
%   % =====================================================================
% INSERT FOR PAPER IV  (Arithmetic Expression Geometry IV)
%
% Source file:  paper-4/sections/04b-transport-density.tex
% Integrated 2026-10-20 as the last subsection of Section 4, immediately after
%   Subsection 4.7 ("Non-exact projective transport and cross-ratio holonomy")
%   and its Example 4.13, via
%   % =====================================================================
% INSERT FOR PAPER IV  (Arithmetic Expression Geometry IV)
%
% Source file:  paper-4/sections/04b-transport-density.tex
% Integrated 2026-10-20 as the last subsection of Section 4, immediately after
%   Subsection 4.7 ("Non-exact projective transport and cross-ratio holonomy")
%   and its Example 4.13, via
%   \input{sections/04b-transport-density}  at the end of
%   paper-4/sections/04-condensation-and-continuation.tex.
%
% Labels required from the host paper (all present):
%   def:projective-edge-transport, thm:projective-torsion-solution,
%   rem:reductive-connection, eq:projective-torsion-components,
%   eq:cross-ratio-holonomy, eq:local-projector,
%   ex:f7-cross-ratio-holonomy, ex:f7-counts.
% New labels introduced here (checked unique in the host paper):
%   subsec:transport-density, lem:edge-factor-triviality,
%   cor:nontrivial-edge, prop:transport-density, prop:transport-normalisation,
%   ex:f7-edge-transports, rem:transport-density-boundary.
%
% Claim status:
%   Lemma 4.14, Corollary 4.15, Propositions 4.16-4.17 : proved
%   Example 4.19                                      : computationally verified
%                                                       (exact arithmetic in F_7)
%   (Section 4 of the current build ends at Example 4.13, so the numbering
%   below starts at 4.14; the subsection is 4.8.)
%
% No physical interpretation is attached to any statement below; in
% particular the transport is not identified with any quantity outside the
% projective framework.
% =====================================================================

\subsection{The transport density and its normalisation}
\label{subsec:transport-density}

Theorem~\ref{thm:projective-torsion-solution} establishes two properties of the
incidence-coupled transport of
Definition~\ref{def:projective-edge-transport}: its exact reduction to affine
torsion on the Borel sector, and the existence of a closed-loop observable that
is not forced to telescope.  The present subsection records three further
elementary facts about that transport, all of them consequences of the
definitions and of the finite computations already available.

First, only certain edges can contribute to the multiplicative component.
Second, on an infinitesimal loop that component has a definite two-form density,
which we compare with the curvature of the reductive connection of
Remark~\ref{rem:reductive-connection}.  Third, the witness loop of
Example~\ref{ex:f7-cross-ratio-holonomy} is computed edge by edge, which
separates the two components of the transport on a single explicit example.

Throughout, $e_k:q_{k-1}\to q_k$ is a transverse lifted step between regular
bivaluations, $\Phi_k\in K^{\times}\cong H$ are its frame coordinates, and
\begin{equation}\label{eq:transport-base-factor}
  \lambda_{\mathrm{base}}(e_k)
  :=\frac{1+a_kb_{k-1}}{1+a_{k-1}b_{k-1}}
  =\frac{\varphi_{b_{k-1}}(v_{a_k})}{\varphi_{b_{k-1}}(v_{a_{k-1}})}
\end{equation}
is the base factor of~\eqref{eq:projective-torsion-components}.  Transversality
gives $1+a_{k-1}b_{k-1}\ne0$, so $\lambda_{\mathrm{base}}(e_k)$ is defined and
nonzero.

\begin{lemma}[When an edge has trivial base factor]
\label{lem:edge-factor-triviality}
For every transverse lifted step $e_k$,
\[
  \lambda_{\mathrm{base}}(e_k)=1
  \quad\Longleftrightarrow\quad
  a_k=a_{k-1}\ \text{or}\ b_{k-1}=0 .
\]
\end{lemma}

\begin{proof}
Both fractions in \eqref{eq:transport-base-factor} share the denominator
$1+a_{k-1}b_{k-1}\ne0$.  Their numerators agree exactly when
$a_kb_{k-1}=a_{k-1}b_{k-1}$, that is $b_{k-1}(a_k-a_{k-1})=0$.
\end{proof}

\begin{corollary}[A necessary condition for a non-unit holonomy]
\label{cor:nontrivial-edge}
Let $\tilde\gamma=(\tilde e_1,\dots,\tilde e_m)$ be a closed lifted path with
$\Phi_m=\Phi_0$.  Then
\[
  \Lambda(\tilde\gamma)=\prod_{k=1}^{m}\lambda_{\mathrm{base}}(e_k),
\]
and $\Lambda(\tilde\gamma)\ne1$ implies that at least one edge satisfies
$b_{k-1}\ne0$ and $a_k\ne a_{k-1}$.  In particular, a closed loop traversed
entirely at kernel coordinate $b=0$, or entirely at constant image coordinate
$a_k\equiv a_{k-1}$, has unit multiplicative holonomy.
\end{corollary}

\begin{proof}
The product formula is the telescoping of the factors $\Phi_k/\Phi_{k-1}$ in
\eqref{eq:projective-torsion-components} around the closed loop.  If every edge
has $b_{k-1}=0$ or $a_k=a_{k-1}$, Lemma~\ref{lem:edge-factor-triviality} makes
every factor $1$, so the product is $1$; this proves the contrapositive of the
stated necessary condition.  The two special cases are the ones in which the
condition fails on every edge.
\end{proof}

The converse of the necessary condition in
Corollary~\ref{cor:nontrivial-edge} is false in general: factors with
$b_{k-1}\ne0$ can multiply to $1$ around a longer loop.  What
Theorem~\ref{thm:projective-torsion-solution}(3) establishes is that on the
four-step loop with both coordinates moved the product is exactly the
cross-ratio \eqref{eq:cross-ratio-holonomy}, which is then non-unit.

\paragraph{The infinitesimal density.}
Let $\square$ be the four-step loop of
Theorem~\ref{thm:projective-torsion-solution}(3) with
$a_2=a_1+\varepsilon$ and $b_2=b_1+\eta$.

\begin{proposition}[Log-holonomy density]
\label{prop:transport-density}
For the loop $\square$,
\begin{equation}\label{eq:log-holonomy}
  \log\Lambda(\square)
  =-\frac{\varepsilon\eta}{(1+a_1b_1)(1+a_2b_2)}+O(\varepsilon^2\eta^2)
  =-\iint_{\square}\frac{\mathrm da\wedge\mathrm db}{(1+ab)^{2}}+O(3),
\end{equation}
where the second expression is the integral of the two-form
$(1+ab)^{-2}\mathrm da\wedge\mathrm db$ over the coordinate rectangle bounded by
$\square$.  Consequently the two-form
\begin{equation}\label{eq:transport-density-form}
  \Omega_{\mathrm{tr}}:=-\frac{\mathrm da\wedge\mathrm db}{(1+ab)^{2}}
\end{equation}
is the density of the multiplicative component: to first order in the size of
the loop, $\log\Lambda(\square)$ is its integral.
\end{proposition}

\begin{proof}
By \eqref{eq:cross-ratio-holonomy},
$\Lambda(\square)=1-\frac{(a_2-a_1)(b_2-b_1)}{(1+a_1b_1)(1+a_2b_2)}
=1-\frac{\varepsilon\eta}{(1+a_1b_1)(1+a_2b_2)}$.  Taking logarithms gives the
first expression.  For the second, note that
$(1+ab)^{-2}=1+O(1)$ near $(a_1,b_1)$ so the integral of the two-form over the
rectangle differs from $\varepsilon\eta/((1+a_1b_1)(1+a_2b_2))$ by a term of
order $3$ in the side lengths, which is absorbed in $O(3)$.
\end{proof}

\paragraph{Comparison with the reductive connection.}
Remark~\ref{rem:reductive-connection} records the connection
$\omega_H=\Phi^{-1}\mathrm d\Phi+(1+ab)^{-1}(a\,\mathrm db-b\,\mathrm da)$ on
$G\to G/H$ together with its curvature
\eqref{eq:reductive-connection-curvature}, which is twice the para-K\"ahler
form.  Comparing that curvature with the density
\eqref{eq:transport-density-form} of Proposition~\ref{prop:transport-density}
leaves a factor of two, and the factor is accounted for by a different
connection, not by a different parametrisation of $\Lambda$.

\begin{proposition}[The transport is the holonomy of the halved transverse part]
\label{prop:transport-normalisation}
Let $f=(1+ab)^{-1}$.  Then
\begin{equation}\label{eq:connection-curvature}
  \mathrm d\omega_H
  =\frac{2\,\mathrm da\wedge\mathrm db}{(1+ab)^{2}}
  =-2\,\Omega_{\mathrm{tr}} ,
\end{equation}
so $\Omega_{\mathrm{tr}}$ is the curvature of the connection
\begin{equation}\label{eq:halved-connection}
  \omega^{(1/2)}=\Phi^{-1}\mathrm d\Phi+\tfrac12\,(1+ab)^{-1}
  \bigl(a\,\mathrm db-b\,\mathrm da\bigr),
\end{equation}
whose transverse part is half that of $\omega_H$ and whose frame part is
unchanged.  In particular the transport of
Definition~\ref{def:projective-edge-transport} agrees, to first order in the
size of the loop, with the holonomy of $\omega^{(1/2)}$, and not with the
holonomy of $\omega_H$.
\end{proposition}

\begin{proof}
The first identity is \eqref{eq:reductive-connection-curvature}, computed in
Remark~\ref{rem:reductive-connection}.  The frame term $\Phi^{-1}\mathrm d\Phi$
is closed, so halving the transverse part halves the curvature and leaves the
frame contribution unchanged; the curvature of \eqref{eq:halved-connection} is
therefore $(1+ab)^{-2}\mathrm da\wedge\mathrm db=-\Omega_{\mathrm{tr}}$.  A
connection with curvature $\Omega$ has holonomy $\exp(-\int_\Sigma\Omega)$
around a small positively oriented loop, which reproduces the density
\eqref{eq:log-holonomy}.
\end{proof}

\begin{remark}[Status of the factor two]
\label{rem:transport-density-boundary}
Proposition~\ref{prop:transport-normalisation} is a normalisation statement, not
a correction: the cross-ratio \eqref{eq:cross-ratio-holonomy} and the connection
of Remark~\ref{rem:reductive-connection} are both exact as they stand, and their
densities differ by $-2$.  Equivalently, to first order in the enclosed area the
holonomy of $\omega_H$ is $\Lambda(\square)^{2}$, not $\Lambda(\square)$; the
idempotent transport is attached to the halved connection
\eqref{eq:halved-connection}.  Nothing in the present framework distinguishes the
two normalisations, no invariant is claimed to, and the choice is left to any
application that fixes a gauge potential; Section~\ref{sec:boundaries} records
the point as the one part of the original projective-torsion problem that
remains open.
\end{remark}

\paragraph{The witness, edge by edge.}
The loop of Example~\ref{ex:f7-cross-ratio-holonomy} is short enough for its
transport to be written down completely.

\begin{example}[Transport of the witness loop over $\F_7$]
\label{ex:f7-edge-transports}
Let $K=\F_7$, $q=7$, and let all frame coordinates be $\Phi_k=1$.  For the loop
with vertices $(a_1,b_1)=(0,0)$, $(a_2,b_1)=(1,0)$, $(a_2,b_2)=(1,1)$,
$(a_1,b_2)=(0,1)$, the transports of
Definition~\ref{def:projective-edge-transport} are, with $2^{-1}=4$ in
$\F_7^{\times}$,
\[
  \mathcal T(e_1)=\begin{pmatrix}1&6\\0&1\end{pmatrix},
  \qquad
  \mathcal T(e_2)=\begin{pmatrix}1&0\\0&1\end{pmatrix},
  \qquad
  \mathcal T(e_3)=\begin{pmatrix}4&4\\0&1\end{pmatrix},
  \qquad
  \mathcal T(e_4)=\begin{pmatrix}1&0\\0&1\end{pmatrix},
\]
so that the only nontrivial step is $e_3$, the step at which the kernel
coordinate moves, and $\lambda_{\mathrm{base}}(e_3)=2^{-1}\equiv4$ in
agreement with Lemma~\ref{lem:edge-factor-triviality} and
\eqref{eq:cross-ratio-holonomy}.
Multiplying chronologically, $\mathcal T(\square)=\mathcal T(e_4)\mathcal T(e_3)
\mathcal T(e_2)\mathcal T(e_1)=\begin{pmatrix}4&0\\0&1\end{pmatrix}$, since
$4\cdot6+4=28\equiv0\pmod 7$.  Thus on this loop the additive component cancels
identically and the non-unit part of the transport is purely multiplicative.
\end{example}

\begin{remark}[What one loop does and does not exhibit]
\label{rem:single-loop-components}
The cancellation in Example~\ref{ex:f7-edge-transports} is a property of that
loop and not of the construction: the additive component is a sum of increments
weighted by different denominators, so nothing forces it to survive.  A loop
witnessing both components simultaneously is obtained by changing one vertex of
$\square$ while keeping the frame lifts closed; we do not compute such a loop
here.  What the example does show is that the two components of
$\mathcal T$ are independent data, and that the finite-field calibration is
strong enough to exhibit them separately.
\end{remark}

These three observations --- the edge criterion of
Corollary~\ref{cor:nontrivial-edge}, the density
\eqref{eq:transport-density-form} together with its normalisation, and the
explicit transport of the witness loop --- are all that the present paper
asserts about the infinitesimal behaviour of the projective torsion.  In
particular no identification of $\Lambda(\square)-1$ with a quantity outside the
projective framework is made, and the choice of normalisation recorded in
Remark~\ref{rem:transport-density-boundary} is left to any later application
that fixes it.
  at the end of
%   paper-4/sections/04-condensation-and-continuation.tex.
%
% Labels required from the host paper (all present):
%   def:projective-edge-transport, thm:projective-torsion-solution,
%   rem:reductive-connection, eq:projective-torsion-components,
%   eq:cross-ratio-holonomy, eq:local-projector,
%   ex:f7-cross-ratio-holonomy, ex:f7-counts.
% New labels introduced here (checked unique in the host paper):
%   subsec:transport-density, lem:edge-factor-triviality,
%   cor:nontrivial-edge, prop:transport-density, prop:transport-normalisation,
%   ex:f7-edge-transports, rem:transport-density-boundary.
%
% Claim status:
%   Lemma 4.14, Corollary 4.15, Propositions 4.16-4.17 : proved
%   Example 4.19                                      : computationally verified
%                                                       (exact arithmetic in F_7)
%   (Section 4 of the current build ends at Example 4.13, so the numbering
%   below starts at 4.14; the subsection is 4.8.)
%
% No physical interpretation is attached to any statement below; in
% particular the transport is not identified with any quantity outside the
% projective framework.
% =====================================================================

\subsection{The transport density and its normalisation}
\label{subsec:transport-density}

Theorem~\ref{thm:projective-torsion-solution} establishes two properties of the
incidence-coupled transport of
Definition~\ref{def:projective-edge-transport}: its exact reduction to affine
torsion on the Borel sector, and the existence of a closed-loop observable that
is not forced to telescope.  The present subsection records three further
elementary facts about that transport, all of them consequences of the
definitions and of the finite computations already available.

First, only certain edges can contribute to the multiplicative component.
Second, on an infinitesimal loop that component has a definite two-form density,
which we compare with the curvature of the reductive connection of
Remark~\ref{rem:reductive-connection}.  Third, the witness loop of
Example~\ref{ex:f7-cross-ratio-holonomy} is computed edge by edge, which
separates the two components of the transport on a single explicit example.

Throughout, $e_k:q_{k-1}\to q_k$ is a transverse lifted step between regular
bivaluations, $\Phi_k\in K^{\times}\cong H$ are its frame coordinates, and
\begin{equation}\label{eq:transport-base-factor}
  \lambda_{\mathrm{base}}(e_k)
  :=\frac{1+a_kb_{k-1}}{1+a_{k-1}b_{k-1}}
  =\frac{\varphi_{b_{k-1}}(v_{a_k})}{\varphi_{b_{k-1}}(v_{a_{k-1}})}
\end{equation}
is the base factor of~\eqref{eq:projective-torsion-components}.  Transversality
gives $1+a_{k-1}b_{k-1}\ne0$, so $\lambda_{\mathrm{base}}(e_k)$ is defined and
nonzero.

\begin{lemma}[When an edge has trivial base factor]
\label{lem:edge-factor-triviality}
For every transverse lifted step $e_k$,
\[
  \lambda_{\mathrm{base}}(e_k)=1
  \quad\Longleftrightarrow\quad
  a_k=a_{k-1}\ \text{or}\ b_{k-1}=0 .
\]
\end{lemma}

\begin{proof}
Both fractions in \eqref{eq:transport-base-factor} share the denominator
$1+a_{k-1}b_{k-1}\ne0$.  Their numerators agree exactly when
$a_kb_{k-1}=a_{k-1}b_{k-1}$, that is $b_{k-1}(a_k-a_{k-1})=0$.
\end{proof}

\begin{corollary}[A necessary condition for a non-unit holonomy]
\label{cor:nontrivial-edge}
Let $\tilde\gamma=(\tilde e_1,\dots,\tilde e_m)$ be a closed lifted path with
$\Phi_m=\Phi_0$.  Then
\[
  \Lambda(\tilde\gamma)=\prod_{k=1}^{m}\lambda_{\mathrm{base}}(e_k),
\]
and $\Lambda(\tilde\gamma)\ne1$ implies that at least one edge satisfies
$b_{k-1}\ne0$ and $a_k\ne a_{k-1}$.  In particular, a closed loop traversed
entirely at kernel coordinate $b=0$, or entirely at constant image coordinate
$a_k\equiv a_{k-1}$, has unit multiplicative holonomy.
\end{corollary}

\begin{proof}
The product formula is the telescoping of the factors $\Phi_k/\Phi_{k-1}$ in
\eqref{eq:projective-torsion-components} around the closed loop.  If every edge
has $b_{k-1}=0$ or $a_k=a_{k-1}$, Lemma~\ref{lem:edge-factor-triviality} makes
every factor $1$, so the product is $1$; this proves the contrapositive of the
stated necessary condition.  The two special cases are the ones in which the
condition fails on every edge.
\end{proof}

The converse of the necessary condition in
Corollary~\ref{cor:nontrivial-edge} is false in general: factors with
$b_{k-1}\ne0$ can multiply to $1$ around a longer loop.  What
Theorem~\ref{thm:projective-torsion-solution}(3) establishes is that on the
four-step loop with both coordinates moved the product is exactly the
cross-ratio \eqref{eq:cross-ratio-holonomy}, which is then non-unit.

\paragraph{The infinitesimal density.}
Let $\square$ be the four-step loop of
Theorem~\ref{thm:projective-torsion-solution}(3) with
$a_2=a_1+\varepsilon$ and $b_2=b_1+\eta$.

\begin{proposition}[Log-holonomy density]
\label{prop:transport-density}
For the loop $\square$,
\begin{equation}\label{eq:log-holonomy}
  \log\Lambda(\square)
  =-\frac{\varepsilon\eta}{(1+a_1b_1)(1+a_2b_2)}+O(\varepsilon^2\eta^2)
  =-\iint_{\square}\frac{\mathrm da\wedge\mathrm db}{(1+ab)^{2}}+O(3),
\end{equation}
where the second expression is the integral of the two-form
$(1+ab)^{-2}\mathrm da\wedge\mathrm db$ over the coordinate rectangle bounded by
$\square$.  Consequently the two-form
\begin{equation}\label{eq:transport-density-form}
  \Omega_{\mathrm{tr}}:=-\frac{\mathrm da\wedge\mathrm db}{(1+ab)^{2}}
\end{equation}
is the density of the multiplicative component: to first order in the size of
the loop, $\log\Lambda(\square)$ is its integral.
\end{proposition}

\begin{proof}
By \eqref{eq:cross-ratio-holonomy},
$\Lambda(\square)=1-\frac{(a_2-a_1)(b_2-b_1)}{(1+a_1b_1)(1+a_2b_2)}
=1-\frac{\varepsilon\eta}{(1+a_1b_1)(1+a_2b_2)}$.  Taking logarithms gives the
first expression.  For the second, note that
$(1+ab)^{-2}=1+O(1)$ near $(a_1,b_1)$ so the integral of the two-form over the
rectangle differs from $\varepsilon\eta/((1+a_1b_1)(1+a_2b_2))$ by a term of
order $3$ in the side lengths, which is absorbed in $O(3)$.
\end{proof}

\paragraph{Comparison with the reductive connection.}
Remark~\ref{rem:reductive-connection} records the connection
$\omega_H=\Phi^{-1}\mathrm d\Phi+(1+ab)^{-1}(a\,\mathrm db-b\,\mathrm da)$ on
$G\to G/H$ together with its curvature
\eqref{eq:reductive-connection-curvature}, which is twice the para-K\"ahler
form.  Comparing that curvature with the density
\eqref{eq:transport-density-form} of Proposition~\ref{prop:transport-density}
leaves a factor of two, and the factor is accounted for by a different
connection, not by a different parametrisation of $\Lambda$.

\begin{proposition}[The transport is the holonomy of the halved transverse part]
\label{prop:transport-normalisation}
Let $f=(1+ab)^{-1}$.  Then
\begin{equation}\label{eq:connection-curvature}
  \mathrm d\omega_H
  =\frac{2\,\mathrm da\wedge\mathrm db}{(1+ab)^{2}}
  =-2\,\Omega_{\mathrm{tr}} ,
\end{equation}
so $\Omega_{\mathrm{tr}}$ is the curvature of the connection
\begin{equation}\label{eq:halved-connection}
  \omega^{(1/2)}=\Phi^{-1}\mathrm d\Phi+\tfrac12\,(1+ab)^{-1}
  \bigl(a\,\mathrm db-b\,\mathrm da\bigr),
\end{equation}
whose transverse part is half that of $\omega_H$ and whose frame part is
unchanged.  In particular the transport of
Definition~\ref{def:projective-edge-transport} agrees, to first order in the
size of the loop, with the holonomy of $\omega^{(1/2)}$, and not with the
holonomy of $\omega_H$.
\end{proposition}

\begin{proof}
The first identity is \eqref{eq:reductive-connection-curvature}, computed in
Remark~\ref{rem:reductive-connection}.  The frame term $\Phi^{-1}\mathrm d\Phi$
is closed, so halving the transverse part halves the curvature and leaves the
frame contribution unchanged; the curvature of \eqref{eq:halved-connection} is
therefore $(1+ab)^{-2}\mathrm da\wedge\mathrm db=-\Omega_{\mathrm{tr}}$.  A
connection with curvature $\Omega$ has holonomy $\exp(-\int_\Sigma\Omega)$
around a small positively oriented loop, which reproduces the density
\eqref{eq:log-holonomy}.
\end{proof}

\begin{remark}[Status of the factor two]
\label{rem:transport-density-boundary}
Proposition~\ref{prop:transport-normalisation} is a normalisation statement, not
a correction: the cross-ratio \eqref{eq:cross-ratio-holonomy} and the connection
of Remark~\ref{rem:reductive-connection} are both exact as they stand, and their
densities differ by $-2$.  Equivalently, to first order in the enclosed area the
holonomy of $\omega_H$ is $\Lambda(\square)^{2}$, not $\Lambda(\square)$; the
idempotent transport is attached to the halved connection
\eqref{eq:halved-connection}.  Nothing in the present framework distinguishes the
two normalisations, no invariant is claimed to, and the choice is left to any
application that fixes a gauge potential; Section~\ref{sec:boundaries} records
the point as the one part of the original projective-torsion problem that
remains open.
\end{remark}

\paragraph{The witness, edge by edge.}
The loop of Example~\ref{ex:f7-cross-ratio-holonomy} is short enough for its
transport to be written down completely.

\begin{example}[Transport of the witness loop over $\F_7$]
\label{ex:f7-edge-transports}
Let $K=\F_7$, $q=7$, and let all frame coordinates be $\Phi_k=1$.  For the loop
with vertices $(a_1,b_1)=(0,0)$, $(a_2,b_1)=(1,0)$, $(a_2,b_2)=(1,1)$,
$(a_1,b_2)=(0,1)$, the transports of
Definition~\ref{def:projective-edge-transport} are, with $2^{-1}=4$ in
$\F_7^{\times}$,
\[
  \mathcal T(e_1)=\begin{pmatrix}1&6\\0&1\end{pmatrix},
  \qquad
  \mathcal T(e_2)=\begin{pmatrix}1&0\\0&1\end{pmatrix},
  \qquad
  \mathcal T(e_3)=\begin{pmatrix}4&4\\0&1\end{pmatrix},
  \qquad
  \mathcal T(e_4)=\begin{pmatrix}1&0\\0&1\end{pmatrix},
\]
so that the only nontrivial step is $e_3$, the step at which the kernel
coordinate moves, and $\lambda_{\mathrm{base}}(e_3)=2^{-1}\equiv4$ in
agreement with Lemma~\ref{lem:edge-factor-triviality} and
\eqref{eq:cross-ratio-holonomy}.
Multiplying chronologically, $\mathcal T(\square)=\mathcal T(e_4)\mathcal T(e_3)
\mathcal T(e_2)\mathcal T(e_1)=\begin{pmatrix}4&0\\0&1\end{pmatrix}$, since
$4\cdot6+4=28\equiv0\pmod 7$.  Thus on this loop the additive component cancels
identically and the non-unit part of the transport is purely multiplicative.
\end{example}

\begin{remark}[What one loop does and does not exhibit]
\label{rem:single-loop-components}
The cancellation in Example~\ref{ex:f7-edge-transports} is a property of that
loop and not of the construction: the additive component is a sum of increments
weighted by different denominators, so nothing forces it to survive.  A loop
witnessing both components simultaneously is obtained by changing one vertex of
$\square$ while keeping the frame lifts closed; we do not compute such a loop
here.  What the example does show is that the two components of
$\mathcal T$ are independent data, and that the finite-field calibration is
strong enough to exhibit them separately.
\end{remark}

These three observations --- the edge criterion of
Corollary~\ref{cor:nontrivial-edge}, the density
\eqref{eq:transport-density-form} together with its normalisation, and the
explicit transport of the witness loop --- are all that the present paper
asserts about the infinitesimal behaviour of the projective torsion.  In
particular no identification of $\Lambda(\square)-1$ with a quantity outside the
projective framework is made, and the choice of normalisation recorded in
Remark~\ref{rem:transport-density-boundary} is left to any later application
that fixes it.
  at the end of
%   paper-4/sections/04-condensation-and-continuation.tex.
%
% Labels required from the host paper (all present):
%   def:projective-edge-transport, thm:projective-torsion-solution,
%   rem:reductive-connection, eq:projective-torsion-components,
%   eq:cross-ratio-holonomy, eq:local-projector,
%   ex:f7-cross-ratio-holonomy, ex:f7-counts.
% New labels introduced here (checked unique in the host paper):
%   subsec:transport-density, lem:edge-factor-triviality,
%   cor:nontrivial-edge, prop:transport-density, prop:transport-normalisation,
%   ex:f7-edge-transports, rem:transport-density-boundary.
%
% Claim status:
%   Lemma 4.14, Corollary 4.15, Propositions 4.16-4.17 : proved
%   Example 4.19                                      : computationally verified
%                                                       (exact arithmetic in F_7)
%   (Section 4 of the current build ends at Example 4.13, so the numbering
%   below starts at 4.14; the subsection is 4.8.)
%
% No physical interpretation is attached to any statement below; in
% particular the transport is not identified with any quantity outside the
% projective framework.
% =====================================================================

\subsection{The transport density and its normalisation}
\label{subsec:transport-density}

Theorem~\ref{thm:projective-torsion-solution} establishes two properties of the
incidence-coupled transport of
Definition~\ref{def:projective-edge-transport}: its exact reduction to affine
torsion on the Borel sector, and the existence of a closed-loop observable that
is not forced to telescope.  The present subsection records three further
elementary facts about that transport, all of them consequences of the
definitions and of the finite computations already available.

First, only certain edges can contribute to the multiplicative component.
Second, on an infinitesimal loop that component has a definite two-form density,
which we compare with the curvature of the reductive connection of
Remark~\ref{rem:reductive-connection}.  Third, the witness loop of
Example~\ref{ex:f7-cross-ratio-holonomy} is computed edge by edge, which
separates the two components of the transport on a single explicit example.

Throughout, $e_k:q_{k-1}\to q_k$ is a transverse lifted step between regular
bivaluations, $\Phi_k\in K^{\times}\cong H$ are its frame coordinates, and
\begin{equation}\label{eq:transport-base-factor}
  \lambda_{\mathrm{base}}(e_k)
  :=\frac{1+a_kb_{k-1}}{1+a_{k-1}b_{k-1}}
  =\frac{\varphi_{b_{k-1}}(v_{a_k})}{\varphi_{b_{k-1}}(v_{a_{k-1}})}
\end{equation}
is the base factor of~\eqref{eq:projective-torsion-components}.  Transversality
gives $1+a_{k-1}b_{k-1}\ne0$, so $\lambda_{\mathrm{base}}(e_k)$ is defined and
nonzero.

\begin{lemma}[When an edge has trivial base factor]
\label{lem:edge-factor-triviality}
For every transverse lifted step $e_k$,
\[
  \lambda_{\mathrm{base}}(e_k)=1
  \quad\Longleftrightarrow\quad
  a_k=a_{k-1}\ \text{or}\ b_{k-1}=0 .
\]
\end{lemma}

\begin{proof}
Both fractions in \eqref{eq:transport-base-factor} share the denominator
$1+a_{k-1}b_{k-1}\ne0$.  Their numerators agree exactly when
$a_kb_{k-1}=a_{k-1}b_{k-1}$, that is $b_{k-1}(a_k-a_{k-1})=0$.
\end{proof}

\begin{corollary}[A necessary condition for a non-unit holonomy]
\label{cor:nontrivial-edge}
Let $\tilde\gamma=(\tilde e_1,\dots,\tilde e_m)$ be a closed lifted path with
$\Phi_m=\Phi_0$.  Then
\[
  \Lambda(\tilde\gamma)=\prod_{k=1}^{m}\lambda_{\mathrm{base}}(e_k),
\]
and $\Lambda(\tilde\gamma)\ne1$ implies that at least one edge satisfies
$b_{k-1}\ne0$ and $a_k\ne a_{k-1}$.  In particular, a closed loop traversed
entirely at kernel coordinate $b=0$, or entirely at constant image coordinate
$a_k\equiv a_{k-1}$, has unit multiplicative holonomy.
\end{corollary}

\begin{proof}
The product formula is the telescoping of the factors $\Phi_k/\Phi_{k-1}$ in
\eqref{eq:projective-torsion-components} around the closed loop.  If every edge
has $b_{k-1}=0$ or $a_k=a_{k-1}$, Lemma~\ref{lem:edge-factor-triviality} makes
every factor $1$, so the product is $1$; this proves the contrapositive of the
stated necessary condition.  The two special cases are the ones in which the
condition fails on every edge.
\end{proof}

The converse of the necessary condition in
Corollary~\ref{cor:nontrivial-edge} is false in general: factors with
$b_{k-1}\ne0$ can multiply to $1$ around a longer loop.  What
Theorem~\ref{thm:projective-torsion-solution}(3) establishes is that on the
four-step loop with both coordinates moved the product is exactly the
cross-ratio \eqref{eq:cross-ratio-holonomy}, which is then non-unit.

\paragraph{The infinitesimal density.}
Let $\square$ be the four-step loop of
Theorem~\ref{thm:projective-torsion-solution}(3) with
$a_2=a_1+\varepsilon$ and $b_2=b_1+\eta$.

\begin{proposition}[Log-holonomy density]
\label{prop:transport-density}
For the loop $\square$,
\begin{equation}\label{eq:log-holonomy}
  \log\Lambda(\square)
  =-\frac{\varepsilon\eta}{(1+a_1b_1)(1+a_2b_2)}+O(\varepsilon^2\eta^2)
  =-\iint_{\square}\frac{\mathrm da\wedge\mathrm db}{(1+ab)^{2}}+O(3),
\end{equation}
where the second expression is the integral of the two-form
$(1+ab)^{-2}\mathrm da\wedge\mathrm db$ over the coordinate rectangle bounded by
$\square$.  Consequently the two-form
\begin{equation}\label{eq:transport-density-form}
  \Omega_{\mathrm{tr}}:=-\frac{\mathrm da\wedge\mathrm db}{(1+ab)^{2}}
\end{equation}
is the density of the multiplicative component: to first order in the size of
the loop, $\log\Lambda(\square)$ is its integral.
\end{proposition}

\begin{proof}
By \eqref{eq:cross-ratio-holonomy},
$\Lambda(\square)=1-\frac{(a_2-a_1)(b_2-b_1)}{(1+a_1b_1)(1+a_2b_2)}
=1-\frac{\varepsilon\eta}{(1+a_1b_1)(1+a_2b_2)}$.  Taking logarithms gives the
first expression.  For the second, note that
$(1+ab)^{-2}=1+O(1)$ near $(a_1,b_1)$ so the integral of the two-form over the
rectangle differs from $\varepsilon\eta/((1+a_1b_1)(1+a_2b_2))$ by a term of
order $3$ in the side lengths, which is absorbed in $O(3)$.
\end{proof}

\paragraph{Comparison with the reductive connection.}
Remark~\ref{rem:reductive-connection} records the connection
$\omega_H=\Phi^{-1}\mathrm d\Phi+(1+ab)^{-1}(a\,\mathrm db-b\,\mathrm da)$ on
$G\to G/H$ together with its curvature
\eqref{eq:reductive-connection-curvature}, which is twice the para-K\"ahler
form.  Comparing that curvature with the density
\eqref{eq:transport-density-form} of Proposition~\ref{prop:transport-density}
leaves a factor of two, and the factor is accounted for by a different
connection, not by a different parametrisation of $\Lambda$.

\begin{proposition}[The transport is the holonomy of the halved transverse part]
\label{prop:transport-normalisation}
Let $f=(1+ab)^{-1}$.  Then
\begin{equation}\label{eq:connection-curvature}
  \mathrm d\omega_H
  =\frac{2\,\mathrm da\wedge\mathrm db}{(1+ab)^{2}}
  =-2\,\Omega_{\mathrm{tr}} ,
\end{equation}
so $\Omega_{\mathrm{tr}}$ is the curvature of the connection
\begin{equation}\label{eq:halved-connection}
  \omega^{(1/2)}=\Phi^{-1}\mathrm d\Phi+\tfrac12\,(1+ab)^{-1}
  \bigl(a\,\mathrm db-b\,\mathrm da\bigr),
\end{equation}
whose transverse part is half that of $\omega_H$ and whose frame part is
unchanged.  In particular the transport of
Definition~\ref{def:projective-edge-transport} agrees, to first order in the
size of the loop, with the holonomy of $\omega^{(1/2)}$, and not with the
holonomy of $\omega_H$.
\end{proposition}

\begin{proof}
The first identity is \eqref{eq:reductive-connection-curvature}, computed in
Remark~\ref{rem:reductive-connection}.  The frame term $\Phi^{-1}\mathrm d\Phi$
is closed, so halving the transverse part halves the curvature and leaves the
frame contribution unchanged; the curvature of \eqref{eq:halved-connection} is
therefore $(1+ab)^{-2}\mathrm da\wedge\mathrm db=-\Omega_{\mathrm{tr}}$.  A
connection with curvature $\Omega$ has holonomy $\exp(-\int_\Sigma\Omega)$
around a small positively oriented loop, which reproduces the density
\eqref{eq:log-holonomy}.
\end{proof}

\begin{remark}[Status of the factor two]
\label{rem:transport-density-boundary}
Proposition~\ref{prop:transport-normalisation} is a normalisation statement, not
a correction: the cross-ratio \eqref{eq:cross-ratio-holonomy} and the connection
of Remark~\ref{rem:reductive-connection} are both exact as they stand, and their
densities differ by $-2$.  Equivalently, to first order in the enclosed area the
holonomy of $\omega_H$ is $\Lambda(\square)^{2}$, not $\Lambda(\square)$; the
idempotent transport is attached to the halved connection
\eqref{eq:halved-connection}.  Nothing in the present framework distinguishes the
two normalisations, no invariant is claimed to, and the choice is left to any
application that fixes a gauge potential; Section~\ref{sec:boundaries} records
the point as the one part of the original projective-torsion problem that
remains open.
\end{remark}

\paragraph{The witness, edge by edge.}
The loop of Example~\ref{ex:f7-cross-ratio-holonomy} is short enough for its
transport to be written down completely.

\begin{example}[Transport of the witness loop over $\F_7$]
\label{ex:f7-edge-transports}
Let $K=\F_7$, $q=7$, and let all frame coordinates be $\Phi_k=1$.  For the loop
with vertices $(a_1,b_1)=(0,0)$, $(a_2,b_1)=(1,0)$, $(a_2,b_2)=(1,1)$,
$(a_1,b_2)=(0,1)$, the transports of
Definition~\ref{def:projective-edge-transport} are, with $2^{-1}=4$ in
$\F_7^{\times}$,
\[
  \mathcal T(e_1)=\begin{pmatrix}1&6\\0&1\end{pmatrix},
  \qquad
  \mathcal T(e_2)=\begin{pmatrix}1&0\\0&1\end{pmatrix},
  \qquad
  \mathcal T(e_3)=\begin{pmatrix}4&4\\0&1\end{pmatrix},
  \qquad
  \mathcal T(e_4)=\begin{pmatrix}1&0\\0&1\end{pmatrix},
\]
so that the only nontrivial step is $e_3$, the step at which the kernel
coordinate moves, and $\lambda_{\mathrm{base}}(e_3)=2^{-1}\equiv4$ in
agreement with Lemma~\ref{lem:edge-factor-triviality} and
\eqref{eq:cross-ratio-holonomy}.
Multiplying chronologically, $\mathcal T(\square)=\mathcal T(e_4)\mathcal T(e_3)
\mathcal T(e_2)\mathcal T(e_1)=\begin{pmatrix}4&0\\0&1\end{pmatrix}$, since
$4\cdot6+4=28\equiv0\pmod 7$.  Thus on this loop the additive component cancels
identically and the non-unit part of the transport is purely multiplicative.
\end{example}

\begin{remark}[What one loop does and does not exhibit]
\label{rem:single-loop-components}
The cancellation in Example~\ref{ex:f7-edge-transports} is a property of that
loop and not of the construction: the additive component is a sum of increments
weighted by different denominators, so nothing forces it to survive.  A loop
witnessing both components simultaneously is obtained by changing one vertex of
$\square$ while keeping the frame lifts closed; we do not compute such a loop
here.  What the example does show is that the two components of
$\mathcal T$ are independent data, and that the finite-field calibration is
strong enough to exhibit them separately.
\end{remark}

These three observations --- the edge criterion of
Corollary~\ref{cor:nontrivial-edge}, the density
\eqref{eq:transport-density-form} together with its normalisation, and the
explicit transport of the witness loop --- are all that the present paper
asserts about the infinitesimal behaviour of the projective torsion.  In
particular no identification of $\Lambda(\square)-1$ with a quantity outside the
projective framework is made, and the choice of normalisation recorded in
Remark~\ref{rem:transport-density-boundary} is left to any later application
that fixes it.
  at the end of
%   paper-4/sections/04-condensation-and-continuation.tex.
%
% Labels required from the host paper (all present):
%   def:projective-edge-transport, thm:projective-torsion-solution,
%   rem:reductive-connection, eq:projective-torsion-components,
%   eq:cross-ratio-holonomy, eq:local-projector,
%   ex:f7-cross-ratio-holonomy, ex:f7-counts.
% New labels introduced here (checked unique in the host paper):
%   subsec:transport-density, lem:edge-factor-triviality,
%   cor:nontrivial-edge, prop:transport-density, prop:transport-normalisation,
%   ex:f7-edge-transports, rem:transport-density-boundary.
%
% Claim status:
%   Lemma 4.14, Corollary 4.15, Propositions 4.16-4.17 : proved
%   Example 4.19                                      : computationally verified
%                                                       (exact arithmetic in F_7)
%   (Section 4 of the current build ends at Example 4.13, so the numbering
%   below starts at 4.14; the subsection is 4.8.)
%
% No physical interpretation is attached to any statement below; in
% particular the transport is not identified with any quantity outside the
% projective framework.
% =====================================================================

\subsection{The transport density and its normalisation}
\label{subsec:transport-density}

Theorem~\ref{thm:projective-torsion-solution} establishes two properties of the
incidence-coupled transport of
Definition~\ref{def:projective-edge-transport}: its exact reduction to affine
torsion on the Borel sector, and the existence of a closed-loop observable that
is not forced to telescope.  The present subsection records three further
elementary facts about that transport, all of them consequences of the
definitions and of the finite computations already available.

First, only certain edges can contribute to the multiplicative component.
Second, on an infinitesimal loop that component has a definite two-form density,
which we compare with the curvature of the reductive connection of
Remark~\ref{rem:reductive-connection}.  Third, the witness loop of
Example~\ref{ex:f7-cross-ratio-holonomy} is computed edge by edge, which
separates the two components of the transport on a single explicit example.

Throughout, $e_k:q_{k-1}\to q_k$ is a transverse lifted step between regular
bivaluations, $\Phi_k\in K^{\times}\cong H$ are its frame coordinates, and
\begin{equation}\label{eq:transport-base-factor}
  \lambda_{\mathrm{base}}(e_k)
  :=\frac{1+a_kb_{k-1}}{1+a_{k-1}b_{k-1}}
  =\frac{\varphi_{b_{k-1}}(v_{a_k})}{\varphi_{b_{k-1}}(v_{a_{k-1}})}
\end{equation}
is the base factor of~\eqref{eq:projective-torsion-components}.  Transversality
gives $1+a_{k-1}b_{k-1}\ne0$, so $\lambda_{\mathrm{base}}(e_k)$ is defined and
nonzero.

\begin{lemma}[When an edge has trivial base factor]
\label{lem:edge-factor-triviality}
For every transverse lifted step $e_k$,
\[
  \lambda_{\mathrm{base}}(e_k)=1
  \quad\Longleftrightarrow\quad
  a_k=a_{k-1}\ \text{or}\ b_{k-1}=0 .
\]
\end{lemma}

\begin{proof}
Both fractions in \eqref{eq:transport-base-factor} share the denominator
$1+a_{k-1}b_{k-1}\ne0$.  Their numerators agree exactly when
$a_kb_{k-1}=a_{k-1}b_{k-1}$, that is $b_{k-1}(a_k-a_{k-1})=0$.
\end{proof}

\begin{corollary}[A necessary condition for a non-unit holonomy]
\label{cor:nontrivial-edge}
Let $\tilde\gamma=(\tilde e_1,\dots,\tilde e_m)$ be a closed lifted path with
$\Phi_m=\Phi_0$.  Then
\[
  \Lambda(\tilde\gamma)=\prod_{k=1}^{m}\lambda_{\mathrm{base}}(e_k),
\]
and $\Lambda(\tilde\gamma)\ne1$ implies that at least one edge satisfies
$b_{k-1}\ne0$ and $a_k\ne a_{k-1}$.  In particular, a closed loop traversed
entirely at kernel coordinate $b=0$, or entirely at constant image coordinate
$a_k\equiv a_{k-1}$, has unit multiplicative holonomy.
\end{corollary}

\begin{proof}
The product formula is the telescoping of the factors $\Phi_k/\Phi_{k-1}$ in
\eqref{eq:projective-torsion-components} around the closed loop.  If every edge
has $b_{k-1}=0$ or $a_k=a_{k-1}$, Lemma~\ref{lem:edge-factor-triviality} makes
every factor $1$, so the product is $1$; this proves the contrapositive of the
stated necessary condition.  The two special cases are the ones in which the
condition fails on every edge.
\end{proof}

The converse of the necessary condition in
Corollary~\ref{cor:nontrivial-edge} is false in general: factors with
$b_{k-1}\ne0$ can multiply to $1$ around a longer loop.  What
Theorem~\ref{thm:projective-torsion-solution}(3) establishes is that on the
four-step loop with both coordinates moved the product is exactly the
cross-ratio \eqref{eq:cross-ratio-holonomy}, which is then non-unit.

\paragraph{The infinitesimal density.}
Let $\square$ be the four-step loop of
Theorem~\ref{thm:projective-torsion-solution}(3) with
$a_2=a_1+\varepsilon$ and $b_2=b_1+\eta$.

\begin{proposition}[Log-holonomy density]
\label{prop:transport-density}
For the loop $\square$,
\begin{equation}\label{eq:log-holonomy}
  \log\Lambda(\square)
  =-\frac{\varepsilon\eta}{(1+a_1b_1)(1+a_2b_2)}+O(\varepsilon^2\eta^2)
  =-\iint_{\square}\frac{\mathrm da\wedge\mathrm db}{(1+ab)^{2}}+O(3),
\end{equation}
where the second expression is the integral of the two-form
$(1+ab)^{-2}\mathrm da\wedge\mathrm db$ over the coordinate rectangle bounded by
$\square$.  Consequently the two-form
\begin{equation}\label{eq:transport-density-form}
  \Omega_{\mathrm{tr}}:=-\frac{\mathrm da\wedge\mathrm db}{(1+ab)^{2}}
\end{equation}
is the density of the multiplicative component: to first order in the size of
the loop, $\log\Lambda(\square)$ is its integral.
\end{proposition}

\begin{proof}
By \eqref{eq:cross-ratio-holonomy},
$\Lambda(\square)=1-\frac{(a_2-a_1)(b_2-b_1)}{(1+a_1b_1)(1+a_2b_2)}
=1-\frac{\varepsilon\eta}{(1+a_1b_1)(1+a_2b_2)}$.  Taking logarithms gives the
first expression.  For the second, note that
$(1+ab)^{-2}=1+O(1)$ near $(a_1,b_1)$ so the integral of the two-form over the
rectangle differs from $\varepsilon\eta/((1+a_1b_1)(1+a_2b_2))$ by a term of
order $3$ in the side lengths, which is absorbed in $O(3)$.
\end{proof}

\paragraph{Comparison with the reductive connection.}
Remark~\ref{rem:reductive-connection} records the connection
$\omega_H=\Phi^{-1}\mathrm d\Phi+(1+ab)^{-1}(a\,\mathrm db-b\,\mathrm da)$ on
$G\to G/H$ together with its curvature
\eqref{eq:reductive-connection-curvature}, which is twice the para-K\"ahler
form.  Comparing that curvature with the density
\eqref{eq:transport-density-form} of Proposition~\ref{prop:transport-density}
leaves a factor of two, and the factor is accounted for by a different
connection, not by a different parametrisation of $\Lambda$.

\begin{proposition}[The transport is the holonomy of the halved transverse part]
\label{prop:transport-normalisation}
Let $f=(1+ab)^{-1}$.  Then
\begin{equation}\label{eq:connection-curvature}
  \mathrm d\omega_H
  =\frac{2\,\mathrm da\wedge\mathrm db}{(1+ab)^{2}}
  =-2\,\Omega_{\mathrm{tr}} ,
\end{equation}
so $\Omega_{\mathrm{tr}}$ is the curvature of the connection
\begin{equation}\label{eq:halved-connection}
  \omega^{(1/2)}=\Phi^{-1}\mathrm d\Phi+\tfrac12\,(1+ab)^{-1}
  \bigl(a\,\mathrm db-b\,\mathrm da\bigr),
\end{equation}
whose transverse part is half that of $\omega_H$ and whose frame part is
unchanged.  In particular the transport of
Definition~\ref{def:projective-edge-transport} agrees, to first order in the
size of the loop, with the holonomy of $\omega^{(1/2)}$, and not with the
holonomy of $\omega_H$.
\end{proposition}

\begin{proof}
The first identity is \eqref{eq:reductive-connection-curvature}, computed in
Remark~\ref{rem:reductive-connection}.  The frame term $\Phi^{-1}\mathrm d\Phi$
is closed, so halving the transverse part halves the curvature and leaves the
frame contribution unchanged; the curvature of \eqref{eq:halved-connection} is
therefore $(1+ab)^{-2}\mathrm da\wedge\mathrm db=-\Omega_{\mathrm{tr}}$.  A
connection with curvature $\Omega$ has holonomy $\exp(-\int_\Sigma\Omega)$
around a small positively oriented loop, which reproduces the density
\eqref{eq:log-holonomy}.
\end{proof}

\begin{remark}[Status of the factor two]
\label{rem:transport-density-boundary}
Proposition~\ref{prop:transport-normalisation} is a normalisation statement, not
a correction: the cross-ratio \eqref{eq:cross-ratio-holonomy} and the connection
of Remark~\ref{rem:reductive-connection} are both exact as they stand, and their
densities differ by $-2$.  Equivalently, to first order in the enclosed area the
holonomy of $\omega_H$ is $\Lambda(\square)^{2}$, not $\Lambda(\square)$; the
idempotent transport is attached to the halved connection
\eqref{eq:halved-connection}.  Nothing in the present framework distinguishes the
two normalisations, no invariant is claimed to, and the choice is left to any
application that fixes a gauge potential; Section~\ref{sec:boundaries} records
the point as the one part of the original projective-torsion problem that
remains open.
\end{remark}

\paragraph{The witness, edge by edge.}
The loop of Example~\ref{ex:f7-cross-ratio-holonomy} is short enough for its
transport to be written down completely.

\begin{example}[Transport of the witness loop over $\F_7$]
\label{ex:f7-edge-transports}
Let $K=\F_7$, $q=7$, and let all frame coordinates be $\Phi_k=1$.  For the loop
with vertices $(a_1,b_1)=(0,0)$, $(a_2,b_1)=(1,0)$, $(a_2,b_2)=(1,1)$,
$(a_1,b_2)=(0,1)$, the transports of
Definition~\ref{def:projective-edge-transport} are, with $2^{-1}=4$ in
$\F_7^{\times}$,
\[
  \mathcal T(e_1)=\begin{pmatrix}1&6\\0&1\end{pmatrix},
  \qquad
  \mathcal T(e_2)=\begin{pmatrix}1&0\\0&1\end{pmatrix},
  \qquad
  \mathcal T(e_3)=\begin{pmatrix}4&4\\0&1\end{pmatrix},
  \qquad
  \mathcal T(e_4)=\begin{pmatrix}1&0\\0&1\end{pmatrix},
\]
so that the only nontrivial step is $e_3$, the step at which the kernel
coordinate moves, and $\lambda_{\mathrm{base}}(e_3)=2^{-1}\equiv4$ in
agreement with Lemma~\ref{lem:edge-factor-triviality} and
\eqref{eq:cross-ratio-holonomy}.
Multiplying chronologically, $\mathcal T(\square)=\mathcal T(e_4)\mathcal T(e_3)
\mathcal T(e_2)\mathcal T(e_1)=\begin{pmatrix}4&0\\0&1\end{pmatrix}$, since
$4\cdot6+4=28\equiv0\pmod 7$.  Thus on this loop the additive component cancels
identically and the non-unit part of the transport is purely multiplicative.
\end{example}

\begin{remark}[What one loop does and does not exhibit]
\label{rem:single-loop-components}
The cancellation in Example~\ref{ex:f7-edge-transports} is a property of that
loop and not of the construction: the additive component is a sum of increments
weighted by different denominators, so nothing forces it to survive.  A loop
witnessing both components simultaneously is obtained by changing one vertex of
$\square$ while keeping the frame lifts closed; we do not compute such a loop
here.  What the example does show is that the two components of
$\mathcal T$ are independent data, and that the finite-field calibration is
strong enough to exhibit them separately.
\end{remark}

These three observations --- the edge criterion of
Corollary~\ref{cor:nontrivial-edge}, the density
\eqref{eq:transport-density-form} together with its normalisation, and the
explicit transport of the witness loop --- are all that the present paper
asserts about the infinitesimal behaviour of the projective torsion.  In
particular no identification of $\Lambda(\square)-1$ with a quantity outside the
projective framework is made, and the choice of normalisation recorded in
Remark~\ref{rem:transport-density-boundary} is left to any later application
that fixes it.
